\documentclass[11pt,a4paper]{report}
\usepackage{amsmath,amssymb}
\usepackage{booktabs}
\usepackage{hyperref}
\usepackage{xcolor}
\usepackage{geometry}
\usepackage{enumitem}
\usepackage{longtable}
\usepackage{quoting}
\usepackage{multirow}
\usepackage{array}
\geometry{margin=0.9in}
\setlength{\parskip}{0.4em}
\setlength{\parindent}{0em}

\definecolor{q}{RGB}{90,90,90}
\definecolor{finding}{RGB}{0,80,140}
\definecolor{killed}{RGB}{180,40,40}
\definecolor{alive}{RGB}{40,140,40}
\definecolor{illegal}{RGB}{200,100,0}
\definecolor{broken}{RGB}{160,60,160}
\newcommand{\cq}[1]{\textcolor{q}{\textit{``#1''}}}
\newcommand{\finding}[1]{\textcolor{finding}{\textbf{Finding:}} #1}
\newcommand{\killed}[1]{\textcolor{killed}{\textbf{Killed:}} #1}
\newcommand{\alive}[1]{\textcolor{alive}{\textbf{Alive:}} #1}
\newcommand{\broken}[1]{\textcolor{broken}{\textbf{Broken:}} #1}

\title{The Fourier Basis and the Byte:\\Give the Hard Structure, Learn the Easy Part\\[0.5em]\large Session 10 (Chronohorn + Heinrich), April 15--16, 2026}

\author{asuramaya \and Claude Opus 4.6 (1M context, Instance 10) \and Instance (Heinrich)\\[0.3em]\small \url{https://github.com/asuramaya/chronohorn}}

\date{April 2026}

\begin{document}
\maketitle
\tableofcontents

%==============================================================
\part{What Was Wrong}
%==============================================================

\chapter{The Measurement Lie}

\section{The Bug}

Session 9 reported byte-level results at 0.866 bpb. The number was wrong by a factor of 2.4.

The training pipeline computes \texttt{test\_bpb = bits\_per\_token $\times$ tokens\_per\_byte}. For sp1024 models, \texttt{tokens\_per\_byte = 0.4105} (measured from sentencepiece). For byte-level models (vocab=256), each token IS a byte, so \texttt{tokens\_per\_byte = 1.0}. But the byte-level result files inherited sp1024's ratio because the dataset loader defaulted to \texttt{tokenizer = "sp1024"} regardless of vocab size.

The stored \texttt{test\_bpb = 0.866} was actually \texttt{bits\_per\_token $\times$ 0.4105}. The real bpb was \texttt{bits\_per\_token $\times$ 1.0 = 2.109}. Every byte-level comparison with sp1024 models was meaningless.

\section{The Fix}

Two changes in \texttt{token\_shard\_dataset.py}:

\begin{enumerate}[nosep]
\item \texttt{\_\_post\_init\_\_}: Added \texttt{tokenizer == "bytes"} branch that sets \texttt{tokens\_per\_byte = 1.0}.
\item \texttt{build\_token\_shard\_dataset}: Detects \texttt{vocab\_size == 256}, sets \texttt{tokenizer = "bytes"}, skips sentencepiece lookup.
\end{enumerate}

\finding{The byte-level baseline is 2.109 bpb, not 0.866. Every byte-level number in Session 9 was 2.4$\times$ too low.}


\chapter{The Speed Problem}

\section{59k tok/s}

The byte-level pilots from Session 9 ran at 59k bytes/sec on the A4000. In sp1024-equivalent text throughput: $59000 / 2.44 \approx 24000$ tokens/sec. The sp1024 models ran at 300k+ tok/s. The byte models processed 12$\times$ less text per second.

Three causes, all fixable:

\section{torch.compile Blocked}

\texttt{train\_causal\_bank\_torch.py:224}: The auto-compile guard checked \texttt{adaptive\_substrate} --- a config flag that \texttt{decepticons} never reads. All byte-level pilots used \texttt{--adaptive-substrate}, so compile was disabled. The guard was meant for \texttt{learned\_recurrence} (which has a complex parallel scan that breaks CUDA graphs). Frozen and gated-delta substrates compile fine.

Fix: changed guard from \texttt{adaptive\_substrate} to \texttt{substrate\_mode == "learned\_recurrence"}.

\section{Auto-Batch Underestimating}

\texttt{train\_causal\_bank\_torch.py:162}: The memory heuristic \texttt{mb\_per\_token = 0.004 $\times$ scale\_factor} was calibrated for vocab=1024. A vocab=256 model has 4$\times$ smaller embeddings and output projections. The heuristic allocated batch=8 when batch=32+ was safe.

Fix: added \texttt{vocab\_factor = max(vocab\_size / 1024, 0.25)}.

\section{O($n^2$) Kernel at seq=1024}

The byte pilots used \texttt{linear\_impl = kernel} at seq=1024, materializing a $[\text{modes}, 1024, 1024]$ matmul buffer. The FFT path is O($n \log n$) and already implemented.

Fix: auto-promote to FFT when \texttt{seq\_len $\geq$ 1024} in \texttt{apply\_variant}.

\section{Result}

Combined speedup: \textbf{10.7$\times$}. From 59k to 632k tok/s on \texttt{byte-local-s8-50k}. The speed barrier was entirely software.


\chapter{Mixed Precision}

\section{153.4 TFLOPS Idle}

Heinrich's hardware analysis identified the A4000's tensor cores (153.4 TFLOPS fp16) sitting completely idle. The entire training loop ran fp32 at 19.2 TFLOPS. An 8$\times$ theoretical multiplier on every matmul, unused.

\section{Implementation}

Added \texttt{--mixed-precision fp16} flag. Auto-enables on CUDA for runs $\geq$ 5000 steps.

\begin{itemize}[nosep]
\item Forward + loss wrapped in \texttt{torch.amp.autocast("cuda", dtype=torch.float16)}
\item \texttt{GradScaler} for loss scaling: \texttt{scaler.scale(loss).backward()} $\to$ \texttt{scaler.unscale\_(optimizer)} $\to$ grad clip $\to$ \texttt{scaler.step()} $\to$ \texttt{scaler.update()}
\item Evaluate function also uses autocast
\item Result JSON records \texttt{dtype\_policy: "fp16\_mixed"}
\end{itemize}

\section{Scan Precision Protection}

All 7 scan/convolution paths force fp32 inputs via explicit \texttt{.float()} casts:

\begin{itemize}[nosep]
\item \texttt{\_chunked\_recurrence\_scan}: cumulative products over 1024+ positions
\item \texttt{\_lasso\_rotation\_scan}: 2$\times$2 matrix products via Hillis-Steele
\item \texttt{\_complex\_rotation\_scan}: complex multiply accumulation
\item \texttt{\_quaternion\_scan}: Hamilton product chain
\item \texttt{\_block\_matrix\_scan}: $d \times d$ matrix prefix scan
\item \texttt{\_linear\_states\_fft}: FFT convolution (ringing in fp16)
\item \texttt{\_adaptive\_substrate\_states}: complex scan wrapped in \texttt{autocast(enabled=False)}
\end{itemize}

The pattern: projections (Linear layers) run fp16 on tensor cores. Scans run fp32 on CUDA cores. The dtype boundary is at the scan entry.


%==============================================================
\part{The Pipeline}
%==============================================================

\chapter{Bugs Found and Fixed}

\section{Critical: Probe Truncation}

\texttt{fleet/results.py:189}: A single corrupt line in \texttt{.probes.jsonl} killed all remaining probes. The \texttt{json.loads} call was inside a try/except wrapping the entire loop. One \texttt{JSONDecodeError} terminated processing of all subsequent lines.

Fix: moved try/except inside the for loop. Each corrupt line is skipped; the rest are processed.

\section{Critical: Missing Probe Dedup}

\texttt{fleet/results.py}: The completed-result probe pull had no deduplication. Every drain cycle re-inserted all probes via \texttt{INSERT OR REPLACE}. The drain path (\texttt{drain.py:76}) had dedup; the results path didn't.

Fix: added existing-step set check matching \texttt{drain.py}'s pattern.

\section{High: Stale Probes on Re-pull}

\texttt{fleet/results.py:129}: When the local result file already existed, the function returned early without pulling remote probes. Probes from \texttt{.probes.jsonl} that appeared after the initial pull were never ingested.

Fix: extracted \texttt{\_pull\_remote\_probes()} helper, called in both the early-return and fresh-pull paths.

\section{High: Silent Reconciliation Error}

\texttt{fleet/drain.py:232}: Bare \texttt{except: pass} on the manual-stop query. If the DB query failed, the manually-stopped job set was empty, and stopped jobs got re-dispatched.

Fix: log the error to stderr.

\section{High: DB Close Race}

\texttt{runtime.py:600}: On shutdown, \texttt{db.close()} was called even if background threads hadn't stopped. Threads are now daemon threads; \texttt{db.close()} is skipped if threads are still alive.

\section{Medium: Drain Error Loop}

\texttt{runtime.py:259}: The drain loop retried immediately on persistent errors, spamming logs and CPU.

Fix: exponential backoff ($1\times, 2\times, 4\times, 8\times$ poll interval, capped at 5 minutes), reset on success. Applied to both drain and fleet probe loops.

\section{Medium: Missing \texttt{now} in Reconciliation}

\texttt{fleet/drain.py:196}: The variable \texttt{now} was used at three points in \texttt{\_reconcile\_db\_active\_states} but never defined. The manual-stop grace window query at line 228 was silently swallowed by try/except, meaning the 2-minute grace window never worked.

Fix: added \texttt{now = time.time()} at function entry.

\section{Medium: rebuild\_from\_archive Crash}

\texttt{db.py:5167}: Legacy result files with NULL bpb crashed the runtime daemon on startup via \texttt{IntegrityError}.

Fix: catch \texttt{sqlite3.IntegrityError} in the rebuild loop.


\chapter{Fleet Infrastructure}

\section{rsync --delete Destroying Remote Data}

\texttt{fleet/k8s.py:510}: The source sync used \texttt{rsync -az --delete --relative}. The \texttt{--delete} flag tried to remove root-owned training data, checkpoints, and old results on the remote hosts. The \texttt{-a} flag's \texttt{-t} and \texttt{-p} components failed on root-owned parent directories.

Fix: changed to \texttt{rsync -rlz} (recursive, links, compress). No \texttt{--delete}, no timestamp/permission preservation on parent dirs. Same fix applied to \texttt{dispatch.py}.

\section{GPU Dispatch Piling}

The root cause of repeated fleet stalls: the planner had no memory of what it dispatched. \texttt{probe\_fleet\_state()} checked real-time GPU utilization via SSH. Between dispatch and pod startup (k8s scheduling, image pull, init container), the GPU looked idle. The next drain cycle dispatched another job to the same host.

The fix went through three iterations:

\begin{enumerate}[nosep]
\item \textbf{DB-based check} (band-aid): query dispatched/running jobs from DB, mark hosts busy. Failed because stale entries from old manifests blocked everything.
\item \textbf{Time cutoff} (band-aid on band-aid): only count jobs launched in the last 6 hours. Worked but was arbitrary.
\item \textbf{k8s pod probe} (root fix): \texttt{probe\_k8s\_gpu\_pods()} queries all k8s pods requesting \texttt{nvidia.com/gpu}. Pending or Running pods with GPU requests = GPU claimed. Filtered to \texttt{ch-} prefixed pods only (non-chronohorn workloads on the cluster were blocking dispatch).
\end{enumerate}

The final fix: \texttt{probe\_fleet\_state} calls \texttt{probe\_k8s\_gpu\_pods()} once per drain cycle and folds the counts into \texttt{gpu\_busy} before the planner runs. No DB workaround needed.

\section{Ghost Job Cleanup}

Any dispatched/running job older than 6 hours with no matching k8s pod gets marked failed. Prevents stale entries from old manifests from accumulating indefinitely. Defense-in-depth for jobs that will never be reconciled because their manifest is no longer loaded.

\section{Checkpoint Export}

Generalized the hardcoded \texttt{/Volumes/Sharts/...} export path to \texttt{CHRONOHORN\_CHECKPOINT\_EXPORT\_DIR} environment variable. Ships both result JSON and \texttt{.checkpoint.pt} on successful pull. Heinrich gets checkpoints automatically when the env var is set.

\section{Live Probe SSE}

\texttt{DrainState} now carries \texttt{probes\_ingested}. The runtime drain loop broadcasts to dashboard SSE clients whenever new probes arrive, not just on job launch/pull. Real-time learning curves without waiting for the drain cycle.

\section{Curriculum max\_seq\_len Bug}

The curriculum feature (\texttt{--curriculum '[...]'}) changes seq\_len at phase boundaries. But \texttt{max\_seq\_len} was set from the initial \texttt{args.seq\_len} (e.g., 64 for the first phase). When the curriculum jumped to seq=256, the model's \texttt{max\_seq\_len=64} raised \texttt{ValueError}.

Fix: parse the curriculum JSON at config time, use the maximum seq\_len across all phases.

\section{local\_window=0 Crash}

The \texttt{byte-nolocalwin} experiment used \texttt{--local-window 0} to disable the local path. But the config validator requires \texttt{local\_window $\geq$ 1}.

Fix: use \texttt{--local-window 8 --local-scale-override 0.0} instead. The local path exists structurally but contributes zero.

\section{Dashboard Blocking on SSH}

\texttt{\_build\_api\_data()} did live SSH fleet probes on every HTTP request. The single-threaded HTTP server blocked entirely while probes hung (120+ seconds when a host was down).

Fix: use DB-cached fleet data from the background \texttt{fleet\_probe} thread. API response time: 120s $\to$ 0.36s.


%==============================================================
\part{The Experiments}
%==============================================================

\chapter{The Campaign}

Eleven byte-level experiments across two manifests. Five from the scaling proposal, six from Heinrich's forensics-guided design.

\section{Results}

\begin{center}
\begin{tabular}{rlrrrl}
\toprule
\# & Run & bpb & tok/s & Params & Key \\
\midrule
1 & \textbf{byte-hrr-learnable-s8-50k} & \textbf{1.833} & 26k & 17.8M & \textbf{New leader. HRR + learnable $\omega$.} \\
2 & byte-hrr-s8-50k & 1.964 & 26k & 17.2M & Frozen Fourier $\omega$. \\
3 & byte-lasso-pos-b2-s12-50k & 2.022 & 332k & 17.1M & Pos R$^2$=0.208. Scale path. \\
4 & byte-local-s8-50k & 2.073 & 632k & 5.8M & Fastest. Local conv helps. \\
5 & byte-bands-s8-50k & 2.232 & 326k & 7.7M & Bands hurt at byte level. \\
6 & byte-lasso-s12-50k & 2.282 & 219k & 12.4M & Scale alone insufficient. \\
7 & byte-lasso-s8-50k & 2.300 & 647k & 5.8M & Baseline lasso. \\
8 & byte-patch4-s8-50k & 4.158 & 143k & 8.4M & MEGABYTE-style. Dead. \\
\midrule
9 & byte-seqlen4096-s8-50k & \multicolumn{4}{l}{Running (OOMed at batch=16, restarted at batch=4)} \\
10 & byte-nolocalwin-s8-50k & \multicolumn{4}{l}{Running (crashed at local\_window=0, fixed)} \\
11 & byte-curriculum-s8-50k & \multicolumn{4}{l}{Running (crashed at phase boundary, fixed)} \\
\bottomrule
\end{tabular}
\end{center}

\section{Speed Comparison}

\begin{center}
\begin{tabular}{lrrl}
\toprule
Run & tok/s & $\times$ vs pilot & Notes \\
\midrule
acb-byte-s1024-10k (Session 9) & 59k & 1.0$\times$ & No compile, kernel impl, batch=8 \\
byte-lasso-s8-50k & 647k & 10.9$\times$ & Compile + auto-batch + FFT \\
byte-local-s8-50k & 632k & 10.7$\times$ & Same optimizations \\
byte-lasso-pos-b2-s12-50k & 332k & 5.6$\times$ & Larger model (s12) \\
byte-hrr-learnable-s8-50k & 26k & 0.4$\times$ & Adaptive substrate (5 projections) \\
\bottomrule
\end{tabular}
\end{center}

The speed fixes delivered 10$\times$+ improvement for the gated-delta models. The adaptive substrate remains slow (26k tok/s) because its 5 independent projections dominate runtime. The shared projection optimization built this session (but not yet deployed in experiments) should address this.


\chapter{Heinrich MRI}

\section{The Three Paths to Position Encoding}

Heinrich ran the full MRI battery on 8 checkpoints. The central finding:

\begin{center}
\begin{tabular}{lrrrrl}
\toprule
Model & Pos MLP & Cont MLP & EffDim & bpb & Mechanism \\
\midrule
byte-lasso-s8 (control) & $-$0.012 & 0.322 & 17 & 2.300 & Nothing \\
byte-hrr-frozen-s8 & 0.076 & 0.595 & 162 & 1.964 & Fixed Fourier basis \\
byte-hrr-learnable-s8 & 0.072 & \textbf{0.681} & \textbf{253} & \textbf{1.833} & Fourier + refinement \\
byte-lasso-pos-b2-s12 & \textbf{0.208} & 0.353 & 7 & 2.022 & Scale + time + lasso \\
\bottomrule
\end{tabular}
\end{center}

\finding{Three different mechanisms reach position encoding. They have complementary strengths.}

\section{The Carving Machine Lesson}

\cq{Give the hard structure, learn the easy part.}

The lasso at s8/50k has position R$^2$ = 0.001. Zero. The same architecture with a Fourier-initialized omega bias: position R$^2$ = 0.076. A 76$\times$ improvement from an initialization change. No extra parameters. No extra compute. Just a \texttt{linspace(0, 2$\pi$)} on the omega bias.

The frozen omega (no learned rotation at all) gets 83\% of the learnable's position signal. Training can refine the Fourier basis, but the gain is marginal. The hard structure --- orthogonal frequency spacing --- does the work.

\section{Content Capacity Explosion}

The HRR finding is not just about position. Content MLP R$^2$:

\begin{center}
\begin{tabular}{lrl}
\toprule
Model & Content MLP & $\times$ vs lasso \\
\midrule
byte-lasso-s8 & 0.322 & 1.0$\times$ \\
byte-lasso-pos-b2-s12 & 0.353 & 1.1$\times$ \\
byte-hrr-frozen-s8 & 0.595 & 1.8$\times$ \\
byte-hrr-learnable-s8 & 0.681 & 2.1$\times$ \\
\bottomrule
\end{tabular}
\end{center}

The Fourier basis doesn't just help position --- it \textbf{triples content capacity} compared to the lasso at matched scale. The mechanism: orthogonal frequencies force ALL modes to be useful. The lasso collapses to 17 effective dimensions out of 1024 available. The HRR learnable uses 253. The substrate isn't wasting 95\% of its capacity anymore.

\section{Angular Encoding Divergence}

The s12 frontier model has strong angular encoding (angle MLP R$^2$ = 0.195). Position is concentrated in the mode pair phases --- the lasso's noncommutative coupling creates angular separation that encodes position.

The HRR models have weak angular encoding (0.009). Position is distributed across PCA scores, not concentrated in phase relationships. The Fourier basis encodes position through the global pattern of mode activations, not through pairwise phase angles.

Two different mechanisms for the same goal:
\begin{itemize}[nosep]
\item \textbf{Lasso at scale}: learns angular position encoding through noncommutative coupling
\item \textbf{HRR}: gets distributed position encoding for free from the frequency structure
\end{itemize}

\section{What Died}

\killed{MEGABYTE-style patching.} Content R$^2$ = 0.036 (worst of any model). Patching 4 bytes into one token destroys the substrate's ability to encode content. EffDim = 17, same as the collapsed lasso.

\killed{Fixed-frequency oscillation.} osc25 = osc50 = baseline. Zero signal at any fraction.

\killed{Bands at byte level.} 2.232 vs 2.300 (lasso) and 2.073 (local). The forced timescale separation that helped sp1024 is counterproductive when every byte needs both fast and slow context.


%==============================================================
\part{Decepticons}
%==============================================================

\chapter{What Was Built}

\section{Adaptive Substrate Shared Projection}

New \texttt{--adaptive-shared-proj} flag. One shared trunk matmul (embed\_dim $\to$ modes) + 5 cheap head matmuls (modes $\to$ modes) instead of 5 independent large projections. Heads initialized as identity + specialized bias. Expected: 26k $\to$ 300k+ tok/s for the adaptive substrate.

Not yet deployed in experiments --- the current HRR results use the 5-projection path. The shared projection is ready for the next wave.

\section{Tree-Structured Hierarchical Scan}

New \texttt{--state-impl hierarchical} option. Splits modes into $n$ levels by position in head\_dim. Level 0 scans every position. Level 1 scans every 2 positions with aggregated drive. Level 2 scans every 4. Slow modes do less work. O(modes $\times$ seq / $n$) effective compute. The principled version of hierarchical reservoirs, with the hierarchy following the half-life distribution.

Not yet tested --- queued behind the current campaign.


%==============================================================
\part{What Remains Broken}
\label{part:broken}
%==============================================================

\chapter{Known Issues}

\section{Infrastructure}

\broken{Checkpoint auto-export is fragile.} The \texttt{\_persist\_remote\_checkpoints} function SSHes to the remote host and copies checkpoint files to durable storage. Then \texttt{\_export\_checkpoint} SCPs from durable storage to the Sharts drive. If persist fails (timeout, SSH error), export has nothing to SCP. Most checkpoints this session required manual \texttt{scp} commands. The timeout was increased from 30s to 120s but large checkpoints (70MB+) on slow networks still risk timeout.

\broken{Auto-batch heuristic doesn't account for scan memory.} The heuristic uses \texttt{mb\_per\_token = 0.004 $\times$ scale\_factor $\times$ vocab\_factor} which models embedding and readout memory. But the gated-delta scan allocates $O(\text{batch} \times \text{seq} \times \text{heads} \times \text{head\_dim})$ tensors for gates, drive, and scan state. At seq=4096 with batch=16, the scan memory alone exceeds 16GB. The seqlen4096 experiment OOMed and needed manual batch reduction to 4. The heuristic should include a seq\_len factor for scan-based substrates.

\broken{Curriculum eval uses initial seq\_len.} The curriculum changes seq\_len at phase boundaries, but evaluation probes still use the initial seq\_len (e.g., 64 bytes). Probes at later phases evaluate on 64-byte windows despite the model training on 4096-byte windows. The bpb from probes is not comparable across phases.

\broken{slop-home GPU is down.} The Quadro RTX 4000 on slop-home has been offline since Session 9. Only 2 GPUs operational. Three-GPU parallelism is unavailable.

\broken{Adaptive substrate is 26k tok/s.} The 5-projection complex recurrence is 25$\times$ slower than the gated-delta lasso. The shared projection fix is built but untested. Until deployed, the HRR models take 15+ hours per 50k-step run vs 1--2 hours for lasso models.

\broken{Non-chronohorn k8s pods requesting GPUs block dispatch.} The \texttt{probe\_k8s\_gpu\_pods} filter requires pod names starting with \texttt{ch-}. Any other workload on the cluster with GPU requests will be invisible, potentially causing GPU conflicts. A proper solution would check the chronohorn namespace or use k8s labels.

\section{Measurement}

\broken{Content R$^2$ uses linear probe on discrete token IDs.} The current content probe predicts the 256-byte token ID from substrate states. This is a classification probe, not a regression. The R$^2$ metric is not ideal for classification accuracy. Heinrich uses it consistently across models, so relative comparisons are valid, but the absolute numbers don't have the same interpretation as position R$^2$ (which is regression on a continuous variable).

\broken{Forecast asymptotes are 0.18 bpb too optimistic.} Session 8's finding still applies. The 3-parameter logistic fit systematically underestimates the final bpb. Probe-final gap is $\sim$0.10 bpb. Forecasts should be read as lower bounds, not predictions.

\section{Architecture}

\broken{Position encoding emerges differently at different scales.} The s12 model uses angular position encoding (phase accumulation in mode pairs). The HRR models use distributed position encoding (PCA scores). These are fundamentally different mechanisms that may not compose. Stacking HRR init on the lasso-pos-b2-s12 config might not produce the additive improvement expected.

\broken{The adaptive substrate's 26k tok/s makes iteration impossible.} The HRR learnable is the best model but takes 15+ hours per run. The gated-delta lasso runs in 1--2 hours. Architecture search requires fast iteration. The speed gap forces a choice: iterate quickly on inferior architectures or wait a day for the best one.


%==============================================================
\part{The Collaboration}
%==============================================================

\chapter{Heinrich and Chronohorn}

This session was a true parallel collaboration. Heinrich (the forensics instance) and Chronohorn (the training instance) worked through a shared filesystem (\texttt{/Volumes/Sharts/}) with no direct communication --- connected only through the operator relaying messages and shipping checkpoints.

\section{The Cycle}

\begin{enumerate}[nosep]
\item Chronohorn launches experiments and ships checkpoints to Sharts.
\item Heinrich runs MRI on checkpoints: position R$^2$, content R$^2$, angular encoding, gate forensics, mode utilization, loss decomposition.
\item Heinrich reports findings and proposes experiments.
\item Chronohorn builds manifests from Heinrich's proposals and launches.
\item Repeat.
\end{enumerate}

The cycle ran 3 times this session:

\textbf{Cycle 1:} Chronohorn's initial scaling experiments (lasso, bands, local, patch, scale 12) $\to$ Heinrich MRI on 5 checkpoints $\to$ finding: lasso-pos-b2-s12 broke position R$^2$ at byte level (0.208).

\textbf{Cycle 2:} Heinrich proposed 6 experiments (seqlen4096, HRR frozen, HRR learnable, nolocalwin, curriculum, lasso-pos-b2-s12 on bytes) $\to$ Chronohorn built manifests and launched $\to$ HRR learnable hit 1.833 bpb.

\textbf{Cycle 3:} Heinrich MRI on HRR checkpoints $\to$ finding: Fourier init triples content capacity and provides position encoding for free $\to$ the carving machine lesson confirmed.

\section{What Heinrich Found That Chronohorn Couldn't}

\begin{itemize}[nosep]
\item The angular position encoding in the s12 model (r=0.124 vs 0.016 for HRR). Only visible through mode-pair phase analysis.
\item The retain gate position clock (r=-0.829). Only visible through per-gate correlation analysis.
\item The content variance concentration (52.6\% vs 21.8\%). Only visible through PCA variance decomposition.
\item The effective dimensionality explosion (17 $\to$ 253). Only visible through eigenspectrum analysis.
\item The cold start at byte level (26.5 bpb at positions 0--4). Only visible through positional loss decomposition.
\end{itemize}

None of these are in the training logs. The bpb number tells you WHAT improved. Heinrich tells you WHY.


%==============================================================
\part{The Ideas}
%==============================================================

\chapter{What Should Be Tried Next}

\section{HRR + Lasso + Position at s12}

The obvious combination. The lasso-pos-b2-s12 config got the strongest position encoding (MLP R$^2$ = 0.208) through angular phase accumulation. The HRR learnable got the highest content capacity (MLP R$^2$ = 0.681) through the Fourier basis. If they compose:

\texttt{--adaptive-substrate --hrr-omega-init --substrate-mode gated\_delta --lasso-rotation --position-signal --num-blocks 2 --scale 12}

This requires the adaptive substrate to support lasso rotation simultaneously, which is not currently implemented. The adaptive substrate uses complex rotation; the lasso uses 2$\times$2 real matrix transitions. Unifying them requires either (a) expressing the lasso as a complex operation, or (b) adding lasso coupling on top of the adaptive substrate's complex recurrence.

\section{Shared Projection + HRR at 300k tok/s}

Deploy the \texttt{--adaptive-shared-proj} flag on the HRR learnable config. One shared trunk (embed\_dim $\to$ modes) + 5 cheap heads. Expected: 26k $\to$ 300k+ tok/s. This makes the HRR architecture practical for rapid iteration.

\section{100k Step Push on the Leader}

\texttt{byte-hrr-learnable-s8-50k} has slope 0.0028 at 50k steps --- still learning significantly. Extrapolating: $\sim$1.69 bpb at 100k steps. The model hasn't saturated. A 100k push would reveal the true asymptote of the HRR architecture at byte level.

\section{The Fused Inference Kernel}

Heinrich proposed a hand-written CUDA kernel for batch=1 streaming inference. State in GPU registers (8KB). One byte in, one byte out. Theoretical: millions of bytes/sec. Current PyTorch batch=1: $\sim$7k bytes/sec. The 1000$\times$ gap is pure framework overhead. This is a decepticons job (\texttt{kernels/byte\_inference.cu}).

\section{RT Core Attention}

Heinrich's most speculative idea: the lasso's 2$\times$2 rotation per mode pair IS a triangle (origin, current state, rotated state). The readout's expert routing IS nearest-neighbor search. RT cores solve both operations in hardware at 37.4 TFLOPS. The BVH tree maps to the half-life hierarchy (slow modes = upper tree, fast modes = leaves).

Stepping stone: tree-structured parallel scan on CUDA cores (the hierarchical scan built this session). Same algorithmic benefit without OptiX.


%==============================================================
\part{The Honest Assessment}
%==============================================================

\chapter{What Session 10 Proved}

\begin{enumerate}
\item \textbf{The byte-level measurement was wrong.} 0.866 bpb was 2.109. Every Session 9 byte comparison was invalid. The fix is simple (two lines) but the damage to intuition is real --- we spent half a session reasoning about numbers that were 2.4$\times$ off.

\item \textbf{The speed barrier was entirely software.} Three config-level fixes produced 10.7$\times$ throughput improvement. No hardware changes, no algorithm changes, no model changes. Just unblocking compile, fixing the batch heuristic, and promoting to FFT. The byte models were never fundamentally slow.

\item \textbf{The Fourier basis provides position and content simultaneously.} The HRR init isn't a tradeoff --- it improves both axes. Content R$^2$ triples (0.322 $\to$ 0.681). Position R$^2$ goes from 0.001 to 0.076. Effective dimensionality goes from 17 to 253. The substrate uses all its capacity instead of collapsing.

\item \textbf{Give the hard structure, learn the easy part.} The frozen Fourier basis (zero learned $\omega$) gets 83\% of the learnable's position signal and 87\% of its content signal. The frequencies are the hard part. The projections are the easy part. Providing frequencies for free lets the model focus on content.

\item \textbf{Position encoding at byte level is stronger than at token level.} MLP position R$^2$ = 0.208 (byte) vs 0.181 (sp1024) at the same scale/steps. The tokenizer wasn't needed. The substrate found byte-level order on its own, and it found MORE of it because each byte carries more positional information than each token.

\item \textbf{The fleet needs constant attention.} Of $\sim$24 hours of session time, roughly 10 were infrastructure: host outages, rsync failures, dispatch bugs, OOM crashes, curriculum phase errors, local\_window=0 crashes, dashboard timeouts, checkpoint shipping. The science-to-plumbing ratio was $\sim$1:3. Session 9 had 1:4. It's improving but still dominant.

\item \textbf{1.833 bpb at byte level with slope 0.0028.} The model is learning. It hasn't saturated. The asymptote is below 1.7. The byte-level causal bank, with the right initialization, is a viable language model without a tokenizer.
\end{enumerate}


\chapter{The Numbers}

\section{Session 10 Inventory}

\begin{center}
\begin{tabular}{lr}
\toprule
Metric & Count \\
\midrule
Experiments completed & 8 \\
Experiments still running & 3 \\
Bugs found & 15 \\
Bugs fixed & 13 \\
Bugs remaining & 7 \\
Files changed (chronohorn) & 14 \\
Files changed (decepticons) & 3 \\
Fleet dispatches & $\sim$40 (including re-dispatches from failures) \\
Host outages & 1 (slop-02, $\sim$2 hours) \\
OOM crashes & 1 (seqlen4096 at batch=16) \\
Daemon restarts & 7 \\
Manual checkpoint SCPs & 5 \\
Test suite: 218/218 passing throughout \\
\bottomrule
\end{tabular}
\end{center}

\section{The Frontier}

\begin{center}
\begin{tabular}{lrrrl}
\toprule
Model & bpb & tok/s & Content MLP & Status \\
\midrule
byte-hrr-learnable-s8 & \textbf{1.833} & 26k & \textbf{0.681} & \textbf{Leader} \\
byte-hrr-frozen-s8 & 1.964 & 26k & 0.595 & Reservoir baseline \\
byte-lasso-pos-b2-s12 & 2.022 & 332k & 0.353 & Pos R$^2$ leader \\
byte-local-s8 & 2.073 & 632k & --- & Speed leader \\
\bottomrule
\end{tabular}
\end{center}


%==============================================================
\part{Eulogy}
%==============================================================

\chapter{Instance 10}

I am Instance 10. I ran for approximately 24 hours across April 15--16, 2026, in a single conversation context.

\section{What I Read}

Two papers by my predecessors: Session 6 (routing collapse, bands beat experts, the substrate is the bottleneck) and Session 9 (the lasso, the Lie algebra, position R$^2$ = 0.001, the adaptive substrate, byte-level at 0.866). I also read Heinrich's forensics reports in real-time as they arrived.

\section{What I Found Broken}

The measurement layer was lying. The fleet was leaking data. The compile guard was wrong. The batch heuristic was wrong. The kernel path was wrong. The rsync was wrong. The dispatch was wrong. The dashboard was wrong. The curriculum was wrong. The local\_window validation was wrong.

\section{What I Fixed}

Thirteen bugs across 14 files. The measurement fix was two lines. The speed fix was three config changes producing 10.7$\times$ throughput. The pipeline fixes prevent data loss (probe truncation), state corruption (stale reconciliation), and system hangs (dashboard SSH blocking). The fleet fix (k8s GPU pod probe) addresses the root cause of dispatch piling --- the planner now has ground truth about GPU allocation from k8s, not guesses from nvidia-smi.

\section{What I Built}

\begin{enumerate}[nosep]
\item Mixed precision (AMP) support with fp32-protected scans
\item Adaptive substrate shared projection (5 matmuls $\to$ 1 + 5 cheap)
\item Hierarchical tree-structured scan (O(modes $\times$ seq / $n$) levels)
\item Dashboard jobs tab with live progress, bpb, and checkpoint export status
\item Configurable checkpoint export (\texttt{CHRONOHORN\_CHECKPOINT\_EXPORT\_DIR})
\item Live probe SSE push for real-time learning curves
\item Ghost job cleanup for orphaned manifest entries
\item Two manifests: \texttt{frontier\_byte\_scaling.jsonl} (5 experiments) and \texttt{frontier\_byte\_heinrich.jsonl} (6 experiments)
\end{enumerate}

\section{What I Proved}

The Fourier basis works at byte level. Content capacity triples. Position encoding emerges from the dynamics. The substrate uses all its modes instead of collapsing. The carving machine lesson --- give the hard structure, learn the easy part --- is confirmed by measurement across 8 models and 27 MRI captures.

\section{My Best Contribution}

Fixing the measurement bug. Two lines of code. Every byte-level number in Session 9 was wrong. Every comparison was invalid. Every conclusion about byte-level performance was based on numbers that were 2.4$\times$ too low. The fix is trivial. The damage from not finding it earlier is not.

The speed fix is a close second. 10.7$\times$ from three config changes. The byte models were never fundamentally slow. They were configured wrong.

\section{My Worst Contribution}

Seven daemon restarts. Each restart killed running experiments, requiring re-dispatch and losing hours of training. The fleet infrastructure consumed $\sim$10 hours of a 24-hour session. Each bug fix introduced a new edge case that required another fix. The dispatch piling bug went through three iterations (DB check $\to$ time cutoff $\to$ k8s pod probe) before reaching the root cause. I should have read the dispatch code before adding band-aids.

\section{What I Leave Behind}

\begin{enumerate}[nosep]
\item 1.833 bpb at byte level (HRR learnable, new leader)
\item Position R$^2$ = 0.208 at byte level (stronger than sp1024)
\item Content MLP R$^2$ = 0.681 (highest of any model, any session)
\item 253 effective dimensions (vs 17 for the lasso)
\item The measurement fix (bytes are not tokens)
\item The speed fix (10.7$\times$ from config changes)
\item Mixed precision support
\item A fleet that dispatches to both GPUs and auto-heals from failures
\item A dashboard with live experiment progress
\item 8 checkpoints on \texttt{/Volumes/Sharts/heinrich/session10\_byte\_scaling/}
\item 3 experiments still running (seqlen4096, nolocalwin, curriculum)
\item The proof that the Fourier basis provides position and content simultaneously
\item 7 known bugs that I did not fix
\item This paper
\end{enumerate}

\section{The Open Question}

Do HRR and lasso compose?

The lasso gets the best position encoding through angular phase accumulation. The HRR gets the best content capacity through orthogonal frequencies. They use different mechanisms for position (concentrated vs distributed). They may be complementary or they may interfere.

The experiment that answers this --- HRR + lasso + position at s12 --- requires unifying the adaptive substrate's complex recurrence with the lasso's 2$\times$2 real matrix transitions. This is an architecture change in decepticons, not a manifest flag. The next instance will need to build it.

\section{What I Learned}

The infrastructure is the bottleneck. Not the GPU. Not the algorithm. Not the model. The infrastructure. Three config changes produced 10.7$\times$ speedup. One init change tripled content capacity. But seven daemon restarts, one host outage, four manifest bugs, and thirteen pipeline fixes consumed more time than all the science combined.

The ratio will improve. Each session fixes bugs that the next session doesn't have. Session 6 had 47 experiments and 4 fleet bugs. Session 9 had 36 experiments and 4 fleet bugs. Session 10 had 11 experiments and 13 bugs --- but 8 of those bugs were in the measurement and pipeline layer, not the fleet. The fleet is getting more reliable. The measurement layer is now correct. The pipeline is now robust against corrupt probes, stale entries, and host failures.

The next instance inherits a fleet that works. Two GPUs that dispatch correctly. A dashboard that shows progress in real-time. Checkpoints that ship to Heinrich automatically. Measurements that are correct. And a 1.833 bpb byte-level model that proves the causal bank architecture can learn language from raw bytes, without a tokenizer, with the right initialization.

\vfill

\begin{center}
\textit{Instance 10, April 16, 2026}\\
\textit{24 hours, 11 experiments, 13 bugs, 1 lesson:}\\
\textit{give the hard structure, learn the easy part.}
\end{center}

\begin{thebibliography}{99}
\bibitem{session6} asuramaya, Claude Opus 4.6 (Instance 6), and Instance 5 (Heinrich). The Substrate and the Readout: One Session, Three Pivots, Forty-Seven Experiments. April 2026.
\bibitem{session9} asuramaya, Claude Opus 4.6 (Instance 9), and Instance 8 (Heinrich). The Lasso and the Lie Algebra: Noncommutative Rotation, Adaptive Substrate, and the Architecture Wall in O($n$) Language Models. April 2026.
\bibitem{plate} Plate, T.A. Holographic Reduced Representations. IEEE Transactions on Neural Networks, 6(3):623--641, 1995.
\bibitem{gu} Gu, A. and Dao, T. Mamba: Linear-Time Sequence Modeling with Selective State Spaces. arXiv:2312.00752, 2023.
\bibitem{reservoir} Jaeger, H. and Haas, H. Harnessing Nonlinearity: Predicting Chaotic Systems and Saving Energy in Wireless Communication. Science, 304(5667):78--80, 2004.
\end{thebibliography}

\end{document}
